\section{敏感性分析与稳健性验证}

为检验模型与结论对关键参数与假设的稳健性，本节从能量余量、资源规模、通信参数、分区数目四个维度展开敏感性分析，并对全文各问的方法对做交叉验证小结。

\subsection{返航安全余量敏感性（问题一）}

返航安全余量 $\rho$ 直接进入能量上限 $(1-\rho)E^{\rm use}_g$，是影响最大安全载荷与组批可行性的关键参数。问题一第 3 小问已在 $\rho\in\{10\%,15\%,20\%,30\%,40\%\}$ 五档下重算 $q_{\max}$ 与组批方案，结果见表~\ref{tab:q1_sens}。当 $\rho$ 由 10\% 提高到 30\%，各机型 $q_{\max}$ 单调下降（C 型由 78.68\,kg 降至 68.51\,kg），能量受限服务区数增加，最少架次由 18 增至 20；当 $\rho=40\%$ 时部分较远服务区出现单箱亦无法安全往返的不可行情形。这一结果表明：20\% 的基准余量处于"载荷能力与返航安全性"的平衡区间，余量设置对远距离、高地形服务区尤为敏感，也解释了后文问题二、三中 C 型机型承担远距离多站路线的瓶颈成因。

\subsection{机队与电池资源敏感性（问题二）}

问题二的资源可行性对机队构成高度敏感。主方案（NSGA-II，29 架次、0 违约）与对比方法（贪心，16 架次、5 违约）的差异本质上是机队瓶颈——C 型仅 2 架机身、4 组电池却承担 10 个架次——造成的排队效应：对比方法将更多货箱压入少量多停靠路线，一旦 C/B 型机身排队，多个硬时限箱被连带拖延而违约；主方案主动拆分重负载路线、增加架次规避拥堵，以更高能耗换取 0 违约。可见"机队规模（尤其 C 型）"与"硬时限可行"之间存在强耦合，若增加 C 型机身或电池组，两类解的架次与违约差距将同步缩小。

充电方面，两阶段充电模型的快充段（SOC$<$90\%）占等效完全充电时间的 65\%，电池充电周转是 makespan 的主要贡献之一。图~\ref{fig:q2_battery} 显示各电池飞行与充电时长量级相当，说明充电周转与飞行任务几乎并行为瓶颈；缩短完全充电时间 $T_{\rm full}$ 可成比例压缩 makespan，但受题设两阶段充电规则约束，本节不改变该口径。

\subsection{通信链路参数敏感性（问题三）}

问题三的通信可行性由自由空间损耗、DEM 地形遮挡（附加损耗 $L_{\rm obs}$）与双向链路预算共同决定，其中遮挡判定与中继悬停高度是主要敏感参数。本文中继悬停位置离散化为"缺口腿水平连线 5 个分数点 $\times$ 6 档高度（50--300\,m）"的候选网格：悬停高度越高，视线越易越过山体遮挡、直连/中继链路越可靠，但中继爬升能耗随高度增加。图~\ref{fig:q3_relay_hover} 显示中继悬停高度主要落在 50--300\,m 内、空间上分布于 DEM 起伏较大的缺口路段两侧，与遮挡判定一致。主方案与对比方法均实现通信中断 0，说明在中继资源（2 架机、6 个能源组件）充足的条件下，通信保障结论对中继候选网格的细分程度不敏感。

地形遮挡判定采用"相位内任一采样点被遮挡即判定全程需中继"的保守口径，会高估中继需求；若改为按子相位区间精细追踪遮挡时段，中继架次与能耗可进一步下降，但这不改变"通信中断可被完全消除"的结论，仅影响中继资源规模的下界。

\subsection{分区数目敏感性（问题四）}

问题四的分区结果随任务组数 $K$ 变化明显。主方案（完全枚举）下 K=2 与 K=3 的总规模分别为 18 与 21、缺口均为 1（中继无人机），均衡系数由 0.0337 上升至 0.3350；对比方法（局部搜索）下缺口由 3 增至 4。总体上资源配置总规模与资源缺口随 $K$ 增大而上升——分组越多，各组须各自留出峰值冗余，同一批次架次被拆到不同组后无法互相借用资源，总需求随 $K$ 上升。这一"分组代价"由"组间不共享"约束内生决定，在 $K$ 取值更大的扩展场景下将更加显著。

\subsection{跨方法交叉验证小结}

表~\ref{tab:robust} 汇总全文四问的方法对与验证结论。各问"主方案 + 对比方法"的结果方向一致、量级可比：问题一三种方法收敛于同一方案；问题二以 0 违约解为主、刻画"严格可行 vs 少架次节能"；问题三两种方法均实现通信中断 0、量化通信代价；问题四精确枚举为局部搜索提供全局最优上界参照，缺口差异被如实报告。多重交叉验证共同支撑了结论的可信度。

\begin{table}[htbp]
\centering
\caption{全文四问方法对与稳健性验证汇总}
\label{tab:robust}
\small
\begin{tabular}{llp{8.5cm}}
\toprule
\textbf{问题} & \textbf{方法对} & \textbf{验证结论} \\
\midrule
问题一 & 线性加权 / Chebyshev / NSGA-II & 三法收敛于 $N=18$、$E=59.072$\,kWh、$T=32801.7$\,s，支付表显示三目标弱冲突 \\
问题二 & NSGA-II（主）/ 贪心（对比） & 0 违约严格可行 vs 16 架次节能但 5 违约，刻画可行域边界 \\
问题三 & 贪心+中继（主）/ NSGA-II 联合（对比） & 均通信中断 0，量化中继能耗与架次的结构性代价 \\
问题四 & 完全枚举（主）/ 局部搜索（对比） & 精确最优缺口 1，局部搜索缺口 3/4，如实报告启发式差距 \\
\bottomrule
\end{tabular}
\end{table}

综上，本文结论对能量余量、资源规模、通信与分区参数的变化方向均有合理解释，且各问方法对结果自洽，稳健性可接受；其中问题二硬时限可行性与 C 型机队规模、问题四资源缺口与组间不共享约束，是最需关注的敏感因素。
